\notechapter{N-20251112}{Calculus Review I: Limits, Completeness, Asymptotic Scales, and Local Approximation}

\section{Cauchy--Schwarz and Minkowski are error-control mechanisms}

For vectors \(x,y\in\mathbb R^n\), the Cauchy--Schwarz inequality is
\[
  |\langle x,y\rangle|
  \le \|x\|_2\|y\|_2.
\]
If \(y\ne0\), consider the nonnegative quadratic
\[
  0\le
  \left\|x-t y\right\|_2^2
  =\|x\|_2^2-2t\langle x,y\rangle+t^2\|y\|_2^2.
\]
Its discriminant is nonpositive, so
\[
  \langle x,y\rangle^2
  \le \|x\|_2^2\|y\|_2^2.
\]
Equality holds exactly when \(x\) and \(y\) are linearly dependent.

Minkowski's inequality follows:
\[
\begin{aligned}
  \|x+y\|_2^2
  &=\|x\|_2^2+2\langle x,y\rangle+\|y\|_2^2\\
  &\le\|x\|_2^2+2\|x\|_2\|y\|_2+\|y\|_2^2\\
  &=(\|x\|_2+\|y\|_2)^2.
\end{aligned}
\]
Hence
\[
  \boxed{\|x+y\|_2\le\|x\|_2+\|y\|_2.}
\]
For nonzero vectors, equality holds exactly when \(x\) and \(y\) point in the same direction.

These inequalities control accumulated errors.  If \(e_1,\ldots,e_m\) are perturbations, then
\[
  \left\|\sum_{j=1}^m e_j\right\|
  \le\sum_{j=1}^m\|e_j\|.
\]
If an error is represented by an inner product, Cauchy--Schwarz converts it into a product of sizes.  The same two patterns survive in Hilbert and normed spaces.

\section{Limits are statements about eventual membership in neighborhoods}

For \(a\in\mathbb R\), let
\[
  N_\varepsilon(a)=\{x:|x-a|<\varepsilon\}.
\]
A sequence \((x_n)\) converges to \(a\) when
\[
  \forall\varepsilon>0,
  \quad
  x_n\in N_\varepsilon(a)
  \text{ eventually}.
\]
Here ``eventually'' means that there exists \(N\) such that the statement holds for every \(n\ge N\).

For a function \(f\),
\[
  \lim_{x\to x_0}f(x)=L
\]
means that for every neighborhood \(V\) of \(L\), there is a punctured neighborhood \(U\) of \(x_0\) such that
\[
  x\in U\quad\Longrightarrow\quad f(x)\in V.
\]
In epsilon--delta notation,
\[
  0<|x-x_0|<\delta
  \quad\Longrightarrow\quad
  |f(x)-L|<\varepsilon.
\]

The neighborhood families form filters: intersections of finitely many neighborhoods are neighborhoods, and larger sets remain neighborhoods.  This is why limits compose.  If
\[
  f(x)\to L
  \quad\text{and}\quad
  g(y)\to M\text{ as }y\to L,
\]
then the preimage under \(g\) of a neighborhood of \(M\) contains a neighborhood of \(L\), and the preimage under \(f\) of that neighborhood contains a punctured neighborhood of \(x_0\).

\section{Sequence limits are unique and respect algebraic operations}

Suppose \(x_n\to a\) and \(x_n\to b\).  If \(a\ne b\), choose
\[
  \varepsilon=\frac{|a-b|}{3}.
\]
For large \(n\),
\[
  |a-b|
  \le|a-x_n|+|x_n-b|
  <2\varepsilon
  <|a-b|,
\]
a contradiction.  Thus limits are unique.

If \(x_n\to a\) and \(y_n\to b\), then
\[
  x_n+y_n\to a+b,
  \qquad
  x_ny_n\to ab.
\]
For the product,
\[
\begin{aligned}
  |x_ny_n-ab|
  &\le |x_n||y_n-b|+|b||x_n-a|.
\end{aligned}
\]
A convergent sequence is bounded, so \(|x_n|\le M\) eventually, and both terms tend to zero.  If \(b\ne0\), then \(y_n\ne0\) eventually and
\[
  \frac{x_n}{y_n}\to\frac ab.
\]

The definition of \(x_n\to+\infty\) is also eventual:
\[
  \forall M>0,
  \quad
  x_n>M\text{ eventually}.
\]
It is not convergence to a real number, but it becomes ordinary convergence after adjoining points at infinity or applying a compactifying coordinate such as \(x\mapsto\arctan x\).

\section{Heine's criterion converts function limits into sequence limits}

\begin{theorem}[Heine criterion]
Let \(f\) be defined on a punctured neighborhood of \(x_0\).  Then
\[
  \lim_{x\to x_0}f(x)=L
\]
iff for every sequence \(x_n\to x_0\) with \(x_n\ne x_0\), one has
\[
  f(x_n)\to L.
\]
\end{theorem}

\begin{proof}
Assume the epsilon--delta limit.  Given \(\varepsilon>0\), choose \(\delta>0\) so that
\[
  0<|x-x_0|<\delta
  \Longrightarrow
  |f(x)-L|<\varepsilon.
\]
Since \(x_n\to x_0\), the hypothesis on \(x_n\) holds eventually, so \(f(x_n)\to L\).

Conversely, suppose the epsilon--delta condition fails.  Then there exists \(\varepsilon_0>0\) such that for every \(n\) one can choose \(x_n\) with
\[
  0<|x_n-x_0|<\frac1n,
  \qquad
  |f(x_n)-L|\ge\varepsilon_0.
\]
Then \(x_n\to x_0\), but \(f(x_n)\not\to L\), contradicting the sequential assumption.
\end{proof}

To disprove a limit, construct incompatible approaches.  For
\[
  f(x)=\cos\frac1x,
\]
take
\[
  x_n=\frac1{2\pi n},
  \qquad
  y_n=\frac1{(2n+1)\pi}.
\]
Both tend to zero, but
\[
  f(x_n)=1,
  \qquad
  f(y_n)=-1.
\]
Therefore \(\lim_{x\to0}f(x)\) does not exist.

\section{Asymptotic notation compares scales, not fixed small numbers}

Let \(g(x)>0\) near a limiting point.  The notation
\[
  f=O(g)
\]
means that there are constants \(C>0\) and a neighborhood on which
\[
  |f(x)|\le Cg(x).
\]
The notation
\[
  f=o(g)
\]
means
\[
  \frac{f(x)}{g(x)}\to0.
\]
We write
\[
  f=\Theta(g)
\]
when both \(f=O(g)\) and \(g=O(f)\), and
\[
  f\sim g
\]
when
\[
  \frac{f(x)}{g(x)}\to1.
\]

The implications are
\[
  f\sim g
  \Longrightarrow
  f=\Theta(g),
  \qquad
  f=o(g)
  \Longrightarrow
  f=O(g),
\]
but the converses fail.

Equivalent infinitesimals may be substituted safely in products and quotients when the remaining factors have controlled limits.  For example,
\[
  \sin x\sim x,
  \qquad
  1-\cos x\sim\frac{x^2}{2}
  \quad(x\to0).
\]
Hence
\[
  \lim_{x\to0}
  \frac{1-\cos x}{x\sin x}
  =\lim_{x\to0}
  \frac{x^2/2}{x^2}
  =\frac12.
\]

Substitution is not valid inside arbitrary sums where cancellation may change the leading order.  Although
\[
  \sqrt{1+x}\sim1,
\]
one cannot replace it by \(1\) in
\[
  \sqrt{1+x}-1.
\]
The difference is of smaller scale:
\[
  \sqrt{1+x}-1
  =\frac{x}{\sqrt{1+x}+1}
  \sim\frac x2.
\]

\section{Squeeze and monotone convergence use order structure}

The squeeze theorem says that if
\[
  a_n\le b_n\le c_n
\]
eventually and
\[
  a_n\to L,
  \qquad
  c_n\to L,
\]
then \(b_n\to L\).  Indeed,
\[
  |b_n-L|
  \le\max\{|a_n-L|,|c_n-L|\}
\]
when \(L\) lies between the bounding terms.

A monotone bounded sequence converges because \(\mathbb R\) has suprema.

\begin{theorem}[Monotone convergence]
If \((x_n)\) is increasing and bounded above, then
\[
  x_n\to\sup_n x_n.
\]
\end{theorem}

\begin{proof}
Let
\[
  L=\sup\{x_n:n\in\mathbb N\}.
\]
Given \(\varepsilon>0\), the number \(L-\varepsilon\) is not an upper bound.  Hence some \(N\) satisfies
\[
  x_N>L-\varepsilon.
\]
For \(n\ge N\), monotonicity gives
\[
  L-\varepsilon<x_N\le x_n\le L.
\]
Thus \(|x_n-L|<\varepsilon\).
\end{proof}

The proof exposes the completeness input: without the least-upper-bound property, the candidate \(L\) may be missing from the number system.

\section{The Cauchy criterion characterizes convergence in complete spaces}

A sequence is Cauchy if
\[
  \forall\varepsilon>0,
  \quad
  |x_m-x_n|<\varepsilon
  \text{ for all sufficiently large }m,n.
\]
Every convergent sequence is Cauchy by the triangle inequality.

For the converse in \(\mathbb R\), first note that a Cauchy sequence is bounded.  Choose \(N\) so that
\[
  |x_n-x_N|<1
  \quad(n\ge N).
\]
The tail lies in a bounded interval, and the finitely many earlier terms are also bounded.

By Bolzano--Weierstrass, a bounded sequence has a convergent subsequence
\[
  x_{n_k}\to L.
\]
Given \(\varepsilon>0\), choose \(K\) so that
\[
  |x_m-x_n|<\frac\varepsilon2
  \quad(m,n\ge K)
\]
and choose \(k\) with \(n_k\ge K\) and
\[
  |x_{n_k}-L|<\frac\varepsilon2.
\]
Then for every \(n\ge K\),
\[
  |x_n-L|
  \le|x_n-x_{n_k}|+|x_{n_k}-L|
  <\varepsilon.
\]
Thus \(x_n\to L\).

The same statement fails in \(\mathbb Q\).  Let
\[
  q_n=\frac{\lfloor10^n\sqrt2\rfloor}{10^n}.
\]
Then \((q_n)\) is Cauchy in the rational metric, but it has no limit in \(\mathbb Q\).  Completion adds the missing limit \(\sqrt2\).

A Banach space is exactly a normed vector space in which every Cauchy sequence converges.  Completeness is therefore not a technical appendix; it is what turns internal error control into the existence of a limit.

\section{Continuity at each point is weaker than uniform continuity}

A function \(f:X\to Y\) between metric spaces is uniformly continuous if
\[
  \forall\varepsilon>0,
  \quad
  \exists\delta>0,
  \quad
  d_X(x,y)<\delta
  \Longrightarrow
  d_Y(f(x),f(y))<\varepsilon
\]
for all \(x,y\in X\).  The same \(\delta\) must work everywhere.

The function \(f(x)=x^2\) is continuous on \(\mathbb R\) but not uniformly continuous.  Set
\[
  x_n=n,
  \qquad
  y_n=n+\frac1n.
\]
Then
\[
  |x_n-y_n|=\frac1n\to0,
\]
but
\[
  |x_n^2-y_n^2|
  =2+\frac1{n^2}
  \not\to0.
\]

On compact domains, continuity becomes uniform.

\begin{theorem}[Heine--Cantor]
A continuous function from a compact metric space to a metric space is uniformly continuous.
\end{theorem}

\begin{proof}
If uniform continuity failed, there would be \(\varepsilon_0>0\) and pairs \(x_n,y_n\) with
\[
  d(x_n,y_n)<\frac1n,
  \qquad
  d(f(x_n),f(y_n))\ge\varepsilon_0.
\]
Compactness gives a convergent subsequence \(x_{n_k}\to x\).  Since
\[
  d(y_{n_k},x)
  \le d(y_{n_k},x_{n_k})+d(x_{n_k},x)
  \to0,
\]
one also has \(y_{n_k}\to x\).  Continuity then forces both image subsequences to converge to \(f(x)\), contradicting their separation by \(\varepsilon_0\).
\end{proof}

Uniform continuity preserves Cauchy sequences and is therefore the natural hypothesis for extending a map to a completion.

\section{A derivative is a first-order approximation with a controlled remainder}

The derivative at \(a\) can be written as
\[
  f(a+h)=f(a)+f'(a)h+r(h),
  \qquad
  \frac{r(h)}{h}\to0.
\]
Equivalently,
\[
  f(a+h)=f(a)+f'(a)h+o(h).
\]
This is stronger than the existence of a slope formula: it says the nonlinear error is negligible compared with the increment.

The differential
\[
  df_a(h)=f'(a)h
\]
is the linear map giving the best first-order approximation.  In several variables, the correct extension is the Fréchet derivative:
\[
  f(a+h)=f(a)+Df(a)h+o(\|h\|).
\]
The derivative is therefore a local linear model, not merely a quotient of small numbers.

The chain rule follows by composing the two local models.  If
\[
  g(a+h)=g(a)+g'(a)h+o(h)
\]
and
\[
  f(g(a)+k)=f(g(a))+f'(g(a))k+o(k),
\]
then \(k=g'(a)h+o(h)=O(h)\), so
\[
  f(g(a+h))
  =f(g(a))+f'(g(a))g'(a)h+o(h).
\]
Thus
\[
  (f\circ g)'(a)=f'(g(a))g'(a).
\]

\section{Taylor's theorem is a finite jet plus a quantified remainder}

If \(f\in C^{n+1}\) near \(a\), Taylor's theorem gives
\[
  f(x)
  =\sum_{k=0}^n
  \frac{f^{(k)}(a)}{k!}(x-a)^k
  +R_n(x).
\]
The Lagrange remainder has the form
\[
  R_n(x)
  =\frac{f^{(n+1)}(\xi)}{(n+1)!}(x-a)^{n+1}
\]
for some \(\xi\) between \(a\) and \(x\).

The integral remainder is
\[
  \boxed{
  R_n(x)
  =\frac1{n!}
  \int_a^x f^{(n+1)}(t)(x-t)^n\,dt.}
\]
It follows by repeated use of the fundamental theorem of calculus.  For example,
\[
\begin{aligned}
  f(x)-f(a)
  &=\int_a^x f'(t_1)\,dt_1\\
  &=f'(a)(x-a)
  +\int_a^x\int_a^{t_1}f''(t_2)\,dt_2\,dt_1,
\end{aligned}
\]
and exchanging the order of integration produces the kernel \((x-t)\).  Iterating gives the general formula.

If
\[
  |f^{(n+1)}(t)|\le M
\]
between \(a\) and \(x\), then
\[
  |R_n(x)|
  \le\frac{M|x-a|^{n+1}}{(n+1)!}.
\]
For \(e^x\) on \([0,1]\), the degree-\(n\) Maclaurin approximation satisfies
\[
  \left|
  e^x-\sum_{k=0}^n\frac{x^k}{k!}
  \right|
  \le\frac{e}{(n+1)!}.
\]
Taylor expansion is therefore an asymptotic statement with a certified error, not permission to replace every smooth function by an infinite series.

\section{Rotation matrices generate trigonometric identities}

Let
\[
  J=
  \begin{pmatrix}
  0&-1\\
  1&0
  \end{pmatrix},
  \qquad
  J^2=-I.
\]
The matrix exponential separates into even and odd powers:
\[
\begin{aligned}
  e^{tJ}
  &=I\cos t+J\sin t\\
  &=
  \begin{pmatrix}
  \cos t&-\sin t\\
  \sin t&\cos t
  \end{pmatrix}.
\end{aligned}
\]
It preserves the Euclidean quadratic form:
\[
  (e^{tJ})^Te^{tJ}=I.
\]

The group law
\[
  e^{(s+t)J}=e^{sJ}e^{tJ}
\]
gives the addition formulas by matrix multiplication:
\[
  \cos(s+t)=\cos s\cos t-\sin s\sin t,
\]
\[
  \sin(s+t)=\sin s\cos t+\cos s\sin t.
\]
Thus the identities are coordinates of the rotation group, not an isolated table to memorize.

\section{Lorentz boosts generate hyperbolic identities}

Let
\[
  K=
  \begin{pmatrix}
  0&1\\
  1&0
  \end{pmatrix},
  \qquad
  K^2=I.
\]
Then
\[
\begin{aligned}
  e^{tK}
  &=I\cosh t+K\sinh t\\
  &=
  \begin{pmatrix}
  \cosh t&\sinh t\\
  \sinh t&\cosh t
  \end{pmatrix}.
\end{aligned}
\]
Let
\[
  \eta=
  \begin{pmatrix}
  1&0\\
  0&-1
  \end{pmatrix}.
\]
A direct multiplication gives
\[
  (e^{tK})^T\eta e^{tK}=\eta.
\]
Therefore the boost preserves
\[
  x^2-y^2.
\]
Applying it to \((1,0)^T\) yields
\[
  \begin{pmatrix}\cosh t\\\sinh t\end{pmatrix},
\]
so invariance gives
\[
  \boxed{\cosh^2t-\sinh^2t=1.}
\]

Again the group law
\[
  e^{(s+t)K}=e^{sK}e^{tK}
\]
produces
\[
  \cosh(s+t)=\cosh s\cosh t+\sinh s\sinh t,
\]
\[
  \sinh(s+t)=\sinh s\cosh t+\cosh s\sinh t.
\]
Circular and hyperbolic functions are parallel because both arise from one-parameter matrix groups.  The sign difference comes from
\[
  J^2=-I,
  \qquad
  K^2=I,
\]
and from the quadratic forms they preserve: Euclidean \(x^2+y^2\) versus Lorentzian \(x^2-y^2\).
